\scanpage{141}
\mspage{90}

en d'autres termes, l'opération $f^*$ sur les préfaisceaux s'interprète, en termes \unsure{géométrique|géométriques}, comme
l'opération $W'\mapsto W'\times_{V'}V$ de changement
de base par $V\unsure{-|\to}V'$, sur les catégories au
dessus de $V'$ qui sont fibrées à fibres discrètes (propriété qui se conserve par changement de base).

Ceci dit, la donnée d'une flèche
$(*)$ équivaut, par la propriété universelle
du changement de base, à celle d'une
flèche (canonique en l'\unsure{occurrence|occurence})
\[
 f_{/F}:V_{/F}\longrightarrow V'_{/f_!(F)}
\]
rendant commutatif le carré
\diagram{\((*)\qquad\begin{tikzcd}[ampersand replacement=\&,row sep=huge,column sep=huge]
 V_{/F} \arrow[r,"{\text{oubli de }F}","{\substack{\text{ou foncteur}\\\text{source}}}"'] \arrow[d,"{f_{/F}}"']
 \& V \arrow[d,"f"]\\
 V'_{/(F'\unsure{{:}{=}|{=}}f_!(F))} \arrow[r,"{\substack{\text{oubli de}\\ F'}}"]
 \& V'
\end{tikzcd}\)}.

Dans le cas qui nous intéresse surtout ici,
$F=e_V$, $V_{/F}\to V$ est un iso., et on
obtient une factorisation canonique de $f$ en

% archive sheet number: 141
